\SourcePage{339}
\SGAsection{8}{爆破簇}
本节给出 SGA 6 第 VII 讲 3.7 和 4.8 中关于\Term{blow-up}{爆破}\Term{scheme}{概形}的 $K$ 群与“\Term{cycle-class-group}{圈类群}”计算的\Term{cohomology}{上同调}变体。假设与记号沿用第 4 节开头（4.1 至 4.3）。在这些条件下，“\Term{fundamental-class}{基本类}”\Term{morphism}{态射}
\[
\gamma:j_*(\mu_{Y'}^{\otimes-1})[-2]\longrightarrow A_{X'}
\]
经\Term{functor}{函子} $\mathrm Rf_*$ \Term{action}{作用}后给出\Term{morphism}{态射}
\[
\mathrm Rf_*(\gamma):i_*\mathrm Rg_*(\mu_{Y'}^{\otimes-1})[-2]\longrightarrow\mathrm Rf_*(A_{X'}).
\]
\begin{SGAresult}{\Term{theorem}{定理}}{8.1}
$D^+(X,A)$ 中的\Term{morphism}{态射}
\[
 A_X\oplus i_*\mathrm Rg_*(\mu_{Y'}^{\otimes-1})[-2]
 \longrightarrow\mathrm Rf_*(A_{X'})\oplus i_*(\mu_Y^{\otimes-d})[-2d],
 \tag{8.1.1}
\]
由矩阵
\[
\begin{pmatrix}\mathrm{adj}&\mathrm Rf_*(\gamma)\\0&i_*(\mathrm{trace})\end{pmatrix}
\]
定义，是\Term{isomorphism}{同构}。
\end{SGAresult}
\begin{proof}
由\Term{stalk-functor}{茎函子}族的\Term{conservative}{保守性}，只需在 $X$ 的每个点 $x$ 上方验证断言。在 $U=X\setminus Y$ 上方断言显然成立。对属于 $Y$ 的 $x$，利用\Term{proper-base-change}{固有基变换定理}，归结为证明以下\Term{lemma}{引理}（取 $r=d-1$）。
\end{proof}
\begin{SGAresult}{\Term{lemma}{引理}}{8.2}
设 $k$ 为\Term{separably-closed}{可分闭}\Term{field}{域}，$r\ge1$ 为整数，$P=\mathbf P_k^r$。则：

a) \Term{canonical-map}{典范映射} $A\longrightarrow H^0(P,A)$ 是\Term{bijection}{双射}。

b) $H^1(P,A)=H^{2r-1}(P,A)=0$。
\SourcePage{340}

c) 对 $0\le p\le2r-2$，以 $\xi=c_1(\mathscr O_P(1))\in H^2(P,\mu)$ 相乘的映射
\[
 H^p(P,\mu_P^{\otimes-1})\longrightarrow H^{p+2}(P,A)
\]
是\Term{isomorphism}{同构}。

d) \Term{trace-morphism}{迹态射} $H^{2r}(P,\mu_P^{\otimes-1})\longrightarrow H^0(k,\mu_k^{\otimes-(r+1)})$ 是\Term{bijection}{双射}。
\end{SGAresult}
此\Term{lemma}{引理}直接由 2.2.2 \Term{pushout}{推出}。

为陈述下一个推论，约定对每个 $X$-\Term{scheme}{概形} $h:T\to X$、$T$ 的每个\Term{closed-subset}{闭子集} $Z$ 和每个 $A_X$-\Term{module}{模}\Term{complex}{复形} $M$，置
\[
 H^*(T,M)=H^*(T,h^*(M)),\qquad H_Z^*(T,M)=H_Z^*(T,h^*(M)).
\]
\begin{SGAresult}{推论}{8.3}
对每个 $A_X$-\Term{module}{模}\Term{complex}{复形} $M$ 及每个 $p\in\mathbf Z$，由矩阵
\[
\begin{pmatrix}f^*&j_*\\0&g_*\end{pmatrix}
\]
定义的映射
\[
 H^p(X,M)\oplus H^{p-2}(Y',M\otimes\mu^{\otimes-1})
 \longrightarrow H^p(X',M)\oplus H^{p-2d}(Y,M\otimes\mu^{\otimes-d})
 \tag{8.3.1}
\]
和
\[
 H_Y^p(X,M)\oplus H^{p-2}(Y',M\otimes\mu^{\otimes-1})
 \longrightarrow H_{Y'}^p(X',M)\oplus H^{p-2d}(Y,M\otimes\mu^{\otimes-d})
 \tag{8.3.2}
\]
都是\Term{bijection}{双射}。
\end{SGAresult}
\begin{proof}
暂以 $u$ 表示\Term{morphism}{态射} (8.1.1)。由\Term{projection-formula}{投影公式}，\Term{arrow}{箭头} (8.3.1) 和 (8.3.2) 分别是\Term{morphism}{态射} $\mathrm R\Gamma(X,u\otimes\mathrm{id}_M)$ 与 $\mathrm R\Gamma_Y(X,u\otimes\mathrm{id}_M)$ 在第 $p$ 个\Term{cohomology}{上同调}\Term{object}{对象}上诱导的\Term{arrow}{箭头}，故得断言。
\end{proof}
\textbf{8.4。}下面将明确写出 (8.3.2) 的逆\Term{isomorphism}{同构}，为此需要几个\Term{lemma}{引理}。
\SourcePage{341}
\begin{SGAresult}{\Term{lemma}{引理}}{8.4.1}
设 $S$ 为\Term{scheme}{概形}，$E$ 为\Term{rank}{秩}恒为 $d$ 的\Term{locally-free}{局部自由} $\mathscr O_S$-\Term{module}{模}，$p:P=\mathbf P_S(E)\to S$ 为相应的\Term{projective-bundle}{射影丛}。以 $F$ 表示由\Term{canonical}{典范}\Term{exact-sequence}{正合列}
\[
0\longrightarrow F\longrightarrow p^*(E)\longrightarrow\mathscr O_P(1)\longrightarrow0
\]
定义的 $\mathscr O_P$-\Term{module}{模}，则有关系
\[
 p_*(c_{d-1}(\check F))=1.
\]
\end{SGAresult}
\begin{proof}
以 $c_i$ 表示 $E$ 的\Term{chern-class}{陈类}，并置 $\xi=c_1(\mathscr O_P(1))$。由等式 $p^*(c(E))=c(F)(1+\xi)$ 得 $c(F)=p^*(c(E))(1+\xi)^{-1}$，从而
\[
 c_{d-1}(F)=\sum_{0\le i\le d-1}(-1)^{d-1-i}c_i\xi^{d-1-i}.
 \tag{8.4.2}
\]
而\Term{homomorphism}{同态} $p_*$ 将\Term{degree}{次数}降低 $2d-2$，且 $p_*(\xi^{d-1})=1$（2.2.2）。对 (8.4.2) 两边施以\Term{homomorphism}{同态} $p_*$，再利用\Term{projection-formula}{投影公式}，即得\Term{lemma}{引理}。
\end{proof}
\begin{SGAresult}{\Term{lemma}{引理}}{8.4.3}
以 $F$ 表示由\Term{canonical}{典范}\Term{exact-sequence}{正合列}
\[
0\longrightarrow F\longrightarrow g^*(N)\longrightarrow\mathscr O_{Y'}(1)\longrightarrow0
\]
定义的\Term{locally-free}{局部自由} $\mathscr O_{Y'}$-\Term{module}{模}。设 $M$ 为在 $Y$ 的某个\Term{neighborhood}{邻域}上\Term{locally-constant}{局部常值}\Term{constructible}{可构造}的 $A_X$-\Term{module}{模}，$p\ge0$ 为整数。对每个 $y\in H^p(Y,M)$，在 $H_{Y'}^{p+2d}(X',M\otimes\mu^{\otimes d})$ 中有关系
\[
 f^*(i_*(y))=j_*(g^*(y)c_{d-1}(\check F)).
\]
\end{SGAresult}
\begin{proof}
置 $z=f^*(i_*(y))-j_*(g^*(y)c_{d-1}(\check F))$。有
\[
 j^*(z)=g^*i^*i_*(y)-j^*j_*(g^*(y)c_{d-1}(\check F)),
\]
由 (4.1) 及\Term{linear}{线性}\Term{lemma}{引理}（SGA 4 第 XVIII 讲）得
\SourcePage{342}
\[
 j^*(z)=g^*(y)[c_d(g^*\check N)-c_{d-1}(\check F)c_1(\mathscr O_{Y'}(-1))]=0.
 \tag{8.4.4}
\]
置 $U=X\setminus Y$、$U'=X'\setminus Y'$，有显然的\Term{commutative-diagram}{交换}\Term{exact-diagram}{正合图}
\[
\begin{tikzcd}[column sep=large]
 H^{p+2d-1}(U')\arrow[r,"\partial"]&H_{Y'}^{p+2d}(X')\arrow[r,"j^*"]&H^{p+2d}(Y')\\
 H^{p+2d-1}(U)\arrow[u,"\wr"]\arrow[r,"\partial"]&H_Y^{p+2d}(X)\arrow[u,"f^*"]\arrow[r,"i^*"]&H^{p+2d}(Y)\arrow[u,"g^*"]
\end{tikzcd}\tag{8.4.5}
\]
其中作了显然的记号简化。由 (8.4.5) 和 (8.4.4) 可知 $z$ 具有\Term{form}{形式} $z=f^*(t)$，其中 $t\in H_Y^{p+2d}(X,M\otimes\mu^{\otimes d})$。由 (4.2)，$t=f_*(z)$，所以只需证明 $f_*(z)=0$。而
\[
 f_*(z)=f_*f^*(i_*(y))-f_*j_*(g^*(y)c_{d-1}(\check F)),
\]
再用 (4.2) 得
\[
 f_*(z)=i_*(y)-i_*g_*(g^*(y)c_{d-1}(\check F)).
\]
利用\Term{morphism}{态射} $g$ 的\Term{projection-formula}{投影公式}及 (8.4.2)，即得结论。
\end{proof}
\begin{SGAresult}{命题}{8.5}
设 $M$ 为在 $Y$ 的某个\Term{neighborhood}{邻域}上\Term{locally-constant}{局部常值}\Term{constructible}{可构造}的 $A_X$-\Term{module}{模}。

a) 以 $\gamma:H_{Y'}^*(X',M\otimes\mu^{\otimes d})\to H^*(Y',M\otimes\mu^{\otimes(d-1)})(-2)$ 表示由关系
\[
 f^*\circ f_*=j_*\circ\gamma+\mathrm{id}\tag{8.5.1}
\]
唯一定义的映射，则
\[
 g_*\circ\gamma=0,\tag{8.5.2}
\]
\SourcePage{343}
且 (8.3.2) 的逆矩阵为
\[
 \begin{pmatrix}f_*&-i_*\\-\gamma&g^*(\,\cdot\,)c_{d-1}(\check F)\end{pmatrix}.
 \tag{8.5.3}
\]

b) 对每个整数 $p\ge0$，序列
\[
0\longrightarrow H^{p-2d}(Y,M(-d))\xrightarrow{\lambda}
 H^{p-2}(Y',M(-1))\oplus H^p(X,M)\xrightarrow{\mu}H^p(X',M)\longrightarrow0,
\]
其中 $\lambda$ 和 $\mu$ 为\Term{arrow}{箭头}
\[
\begin{aligned}
\lambda:x&\longmapsto[g^*(x)c_{d-1}(\check F),-i_*(x)],\\
\mu:(x,y)&\longmapsto j_*(x)+f^*(y),
\end{aligned}
\]
是\Term{exact-sequence}{正合列}。此外，映射 $\lambda$ 有\Term{left-inverse}{左逆}
\[
 \lambda':(x,y)\longmapsto g_*(x).
\]
\end{SGAresult}
\begin{proof}
先证 a)。由 4.2（$f_*f^*=\mathrm{id}$），要证明
\[
\begin{pmatrix}f_*&-i_*\\-\gamma&g^*(\,\cdot\,)c_{d-1}(\check F)\end{pmatrix}
\circ\begin{pmatrix}f^*&j_*\\0&g_*\end{pmatrix}=\mathrm{id},
\]
只需证明
\[
 \gamma j_*+\mathrm{id}=g^*[g_*(\,\cdot\,)c_{d-1}(\check F)]\tag{8.5.4}
\]
和
\[
 \gamma f^*=0.\tag{8.5.5}
\]
由 (8.5.1) 得 $f^*f_*f^*=j_*\gamma f^*+f^*$，由 (4.2) 得 $j_*\gamma f^*=0$；因 $j_*$ 为\Term{isomorphism}{同构}，便得 (8.5.5)。再证 (8.5.4)。由 (8.5.1) 得
\SourcePage{344}
\[
 f^*f_*j_*=j_*\gamma j_*+j_*.
\]
因 $f_*j_*=i_*g_*$，并由 (8.4.3)，得
\[
 j_*g^*(g_*(\,\cdot\,)c_{d-1}(\check F))=j_*\gamma j_*+j_*.
\]
左乘 $j_*$ 的逆，即得关系 (8.5.4)。最后，将这两个矩阵按另一顺序复合并写出它等于\Term{identity-map}{恒等映射}，即得 (8.5.2)。再证 b)。由 8.4.3 和 8.4.1 分别得 $\mu\circ\lambda=0$ 和 $\lambda'\circ\lambda=\mathrm{id}$。另一方面，由 (8.3)，映射 $\mu$ 为\Term{surjective}{满射}。故只需证明 $\operatorname{Ker}(\mu)\subset\operatorname{Im}(\lambda)$。但 $\mu\circ(\lambda\circ\lambda')=0$ 且 $\lambda'\circ(\lambda\circ\lambda')=\mathrm{id}$，所以 $\mu(x)=0$ 蕴含
\[
 (\mu\oplus\lambda')(\lambda\circ\lambda'(x))=(\mu\oplus\lambda')(x),
\]
由 8.3 得 $x=\lambda\circ\lambda'(x)$。
\end{proof}
\textbf{8.6。}现在尝试确定 $X'$ 的\Term{cohomology}{上同调}的\Term{multiplication}{乘法}结构。为此，对每个 $p\in\mathbf Z$ 以及在 $Y$ 的某个\Term{neighborhood}{邻域}上\Term{locally-constant}{局部常值}\Term{constructible}{可构造}的每个 $A_X$-\Term{module}{模} $M$，置
\[
 \nabla^p(X',M)=H^{p-2}(Y',M(-1))\oplus H^p(X,M).
\]
对 $p,q\in\mathbf Z$ 以及在 $Y$ 的某个\Term{neighborhood}{邻域}上\Term{locally-constant}{局部常值}\Term{constructible}{可构造}的两个 $A_X$-\Term{module}{模} $M,N$，定义\Term{composition-law}{合成律}
\[
 \mathrm U:\nabla^p(X',M)\otimes_A\nabla^q(X',N)
 \longrightarrow\nabla^{p+q}(X',M\otimes_A N),
 \tag{8.6.1}
\]
其公式为
\[
 (y'_1,x_1)\mathbin{\mathrm U}(y'_2,x_2)=
 (-y'_1y'_2\xi+y'_1g^*i^*(x_2)+(-1)^{pq}y'_2g^*i^*(x_1),x_1x_2).
 \tag{8.6.2}
\]
此\Term{composition-law}{合成律}显然\Term{anticommutative}{反交换}；一个简单但冗长的计算表明它满足\Term{associative-law}{结合律}。特别地，
\[
 \nabla^*(X',A)=\bigoplus_{p\in\mathbf Z}\nabla^p(X',A)
\]
\SourcePage{345}
也具有\Term{graded-commutative-algebra}{分次交换 $A$-代数}结构（按 $\mathbf Z$ \Term{graded}{分次}的\Term{algebra}{代数}之意义），而
\[
 \nabla^*(X',M)=\bigoplus_{p\in\mathbf Z}\nabla^p(X',M)
\]
具有\Term{graded-ring}{分次环} $\nabla^*(X',A)$ 上的\Term{graded-module}{分次模}结构。
\begin{SGAresult}{命题}{8.6.3}
a) 对在 $Y$ 的某个\Term{neighborhood}{邻域}上\Term{locally-constant}{局部常值}\Term{constructible}{可构造}的每对 $A_X$-\Term{module}{模} $(M,N)$，图
\[
\begin{tikzcd}[column sep=large]
\nabla^*(X',M)\otimes\nabla^*(X',N)\arrow[r,"\mu\otimes\mu"]\arrow[d,"{(8.6.1)}"']&H^*(X',M)\otimes H^*(X',N)\arrow[d,"\cup"]\\
\nabla^*(X',M\otimes N)\arrow[r,"\mu"]&H^*(X',M\otimes N)
\end{tikzcd}
\]
\Term{commutative-diagram}{交换}。特别地，\Term{surjective}{满射}
\[
 \mu_A=\mu:\nabla^*(X',A)\longrightarrow H^*(X',A)
\]
是\Term{ring-homomorphism}{环同态}，而映射
\[
 \mu_M=\mu:\nabla^*(X',M)\longrightarrow H^*(X',M)
\]
是底环可变的\Term{module}{模}\Term{morphism}{态射}。

b) 若通过映射
\[
\begin{aligned}
\nabla^p(X',A)\otimes H^{q-2d}(Y,M(-d))&\longrightarrow H^{p+q-2d}(Y,M(-d)),\\
 (y',x)\otimes y&\longmapsto i^*(x)\cdot y
\end{aligned}
\]
使 $\nabla^*(X',A)$ \Term{action}{作用}于 $H^*(Y,M(-d))(-2d)$，则 8.5 b) 中的映射 $\lambda$ 成为 $\nabla^*(X',A)$-\Term{module}{模}\Term{morphism}{态射}。
\end{SGAresult}
\begin{proof}
先证 a)。用 8.6.2 的记号，需证等式
\[
\begin{aligned}
 {}&j_*(-y'_1y'_2\xi+y'_1g^*i^*(x_2)+(-1)^{pq}y'_2g^*i^*(x_1))+f^*(x_1x_2)\\
 &\qquad=(j_*(y'_1)+f^*(x_1))(j_*(y'_2)+f^*(x_2)).
\end{aligned}
\]
\SourcePage{346}
这直接由下列等式得到：
\[
\begin{aligned}
 f^*(x_1)f^*(x_2)&=f^*(x_1x_2)&&\text{（$f^*$ 的乘法性）},\\
 j_*(-y'_1y'_2\xi)&=j_*(y'_1)j_*(y'_2)&&\text{（4.5）},\\
 j_*(y'g^*i^*(x))&=j_*(y'j^*f^*(x))=j_*(y')f^*(x)&&\text{（\Term{projection-formula}{投影公式}）}.
\end{aligned}
\]
再证 b)。所需的等式 $\lambda(i^*(x)\cdot y)=(y',x)\cdot\lambda(y)$ \Term{resolution}{分解}为以下两式：
\[
\begin{aligned}
 g^*(i^*(x)\cdot y)c_{d-1}(\check F)
 &=[-y'g^*(y)c_{d-1}(\check F)\xi+y'g^*i^*(-i_*(y))\\
 &\qquad+g^*i^*(x)g^*(y)c_{d-1}(\check F)],
\end{aligned}
\]
以及 $-i_*(i^*(x)y)=-x\cdot i_*(y)$。第二式正是\Term{morphism}{态射} $i$ 的\Term{projection-formula}{投影公式}。第一式直接由下列关系得到：
\[
\begin{aligned}
 -\xi c_{d-1}(\check F)&=g^*[c_d(\check N)]&&\text{（8.4.3 的\Term{exact-sequence}{正合列}）},\\
 i^*i_*(y)&=y\,c_d(\check N)&&\text{（4.1）}.
\end{aligned}
\]
\end{proof}
\SourcePage{347}
\SGAsection{9}{爆破簇的 Chow 环及 Chow 环中的自交公式}
本节旨在给出 Mumford 对 4.1 在\Term{chow-ring}{Chow 环}中的对应命题的证明。随后说明如何由此\Term{pushout}{推出}“关键公式”(8.4.3)，进而确定\Term{blow-up}{爆破}\Term{variety}{簇}的\Term{chow-ring}{Chow 环}结构。

依照 3.9 的约定，以 $C^\bullet(X)$ 表示\Term{field}{域} $k$ 上\Term{quasi-projective}{拟射影}\Term{smooth}{光滑}\Term{scheme}{概形} $X$ 的\Term{chow-ring}{Chow 环}。所用的\Term{chern-class}{陈类}理论为 [1] 中的理论。

我们将多次使用以下\Term{lemma}{引理}。
\begin{SGAresult}{\Term{lemma}{引理}}{9.1}
设 $S$ 为\Term{separably-closed}{可分闭}\Term{field}{域} $k$ 上的\Term{quasi-projective}{拟射影}\Term{smooth}{光滑}\Term{scheme}{概形}，$E$ 为\Term{rank}{秩}恒为 $r+1$ 的\Term{locally-free}{局部自由} $\mathscr O_S$-\Term{module}{模}，$p:P=\mathbf P_S(E)\to S$ 为相应的\Term{projective-bundle}{射影丛}。以 $F$ 表示由\Term{canonical}{典范}\Term{exact-sequence}{正合列}
\[
 0\longrightarrow F\longrightarrow p^*(E)\longrightarrow\mathscr O_P(1)\longrightarrow0
 \tag{9.1.0}
\]
定义的 $\mathscr O_P$-\Term{module}{模}，并置 $\xi=c_1(\mathscr O_P(1))$。则：

(i) 对每个 $y\in C^\bullet(S)$，有
\[
 y=p_*(p^*(y)c_r(\check F)).\tag{9.1.1}
\]

(ii) 设 $E=N\oplus\mathscr O_S$，其中 $N$ 为\Term{rank}{秩}为 $r$ 的\Term{locally-free}{局部自由} $\mathscr O_S$-\Term{module}{模}，从而 $P$ 可等同于 $\mathbf V(N)$ 的\Term{projective-completion}{射影完备化} $\widehat N$；以 $s:S\to\mathbf P_S(E)$ 表示 $\mathbf V(N)$ 的\Term{zero-section}{零截面}与\Term{canonical}{典范}\Term{inclusion}{包含映射} $\mathbf V(N)\hookrightarrow\widehat N$ 的复合。则：

(a) 对每个 $y\in C^\bullet(S)$，有关系
\[
 s_*(y)=p^*(y)c_r(\check F)=p^*(y)\left(\sum_{i=0}^r p^*(c_i(\check N))\xi^{r-i}\right),
 \quad\text{在 }C^\bullet(\widehat N)\text{ 中},\tag{9.1.2}
\]
\[
 s^*s_*(y)=y\,c_r(\check N),\quad\text{在 }C^\bullet(S)\text{ 中}.\tag{9.1.3}
\]
\SourcePage{348}

(b) 对每个 $z\in C^\bullet(\widehat N)$，有
\[
 s^*(z)=p_*(z\,c_r(\check F)),\quad\text{在 }C^\bullet(S)\text{ 中}.\tag{9.1.4}
\]
\end{SGAresult}
\begin{proof}
先证 (9.1.1)。由 $p$ 的\Term{projection-formula}{投影公式}，只需证明 $p_*(c_r(\check F))=1$。为此，只须照用 8.4.1 的证明，并利用关系（[CH] 4 (17)）。(9.1.2) 第一项与第三项相等，正是 [1] 的\Term{lemma}{引理} 3。仿照 (8.4.1) 的证明可见
\[
 c_r(\check F)=\sum_{i=0}^r p^*(c_i(\check N))\xi^{r-i},\tag{9.1.2 bis}
\]
这便完成了 (9.1.2) 的证明。对 (9.1.2) 的首尾两项施以\Term{functor}{函子} $s^*$，因 $p\circ s=\mathrm{id}$，可知 (9.1.3) 将由等式
\[
 s^*\left(\sum_{i=0}^{r-1}p^*(c_i(\check N))\xi^{r-i}\right)=0
\]
\Term{pushout}{推出}，更可由等式
\[
 s^*(\xi^i)=0\qquad(i>0)\tag{9.1.5}
\]
\Term{pushout}{推出}。现证 (9.1.5)。以 $\alpha:\mathbf V(N)\hookrightarrow\widehat N$ 表示\Term{canonical}{典范}\Term{inclusion}{包含映射}，显然 $\alpha^*(\mathscr O_P(1))\simeq\mathscr O_{\mathbf V(N)}$，由\Term{chern-class}{陈类}的\Term{functoriality}{函子性}得
\[
 \alpha^*(\xi)=\alpha^*(c_1(\mathscr O_P(1)))=c_1(\mathscr O_{\mathbf V(N)})=0.
\]
由 $s$ 的定义立即得到断言。再证 (9.1.4)。为此使用等式
\[
 \xi c_r(\check F)=c_{r+1}(p^*(\check N)\oplus\mathscr O_P)=0,\tag{9.1.6}
\]
它直接由\Term{exact-sequence}{正合列} (9.1.0) 得到。元素 $z$ 有唯一\Term{resolution}{分解}
\[
 z=\sum_{i=0}^r p^*(a_i)\xi^i,\qquad a_i\in C^\bullet(S)
\]
（[CH] 4 \Term{theorem}{定理} 1）。由 (9.1.5) 得 $s^*(z)=a_0$。另一方面，等式 (9.1.6) 表明 $z\,c_r(\check F)=p^*(a_0)c_r(\check F)$，再由 (9.1.1) 即得结论。
\end{proof}
\SourcePage{349}
现在陈述\Term{chow-ring}{Chow 环}中的“\Term{self-intersection}{自交}公式”。
\begin{SGAresult}{\Term{theorem}{定理}}{9.2（Mumford）}
设 $k$ 为\Term{separably-closed}{可分闭}\Term{field}{域}，$i:Y\hookrightarrow X$ 为\Term{quasi-projective}{拟射影}\Term{smooth}{光滑} $k$-\Term{scheme}{概形}间的\Term{closed-immersion}{闭浸入}，纯\Term{codimension}{余维数}为 $d$。以 $N$ 表示 $i$ 的\Term{normal-sheaf}{法层}，则对每个 $y\in C^\bullet(Y)$，在 $C^\bullet(Y)$ 中有等式
\[
 i^*i_*(y)=y\,c_d(\check N).
\]
\end{SGAresult}
为证明此\Term{theorem}{定理}，需要若干记号与辅助\Term{scheme}{概形}。

借助\Term{morphism}{态射} $u:t\mapsto(t,0)$，将 $X$ \Term{embedding}{嵌入} $Z=X\times_k\mathbf P_k^1$，从而有\Term{cartesian-diagram}{笛卡尔图}
\[
\begin{tikzcd}
\widehat N\arrow[r,hook,"j_1"]\arrow[d,"g_1"']&Z'\arrow[d,"f_1"]\\
Y\arrow[r,hook,"i_1"]&Z
\end{tikzcd}
\]
其中 $Z'$ 表示 $Z$ 沿 $Y$ 的\Term{blow-up}{爆破}，$\widehat N=\mathbf P_Y(N\oplus\mathscr O_Y)$ 为\Term{vector-bundle}{向量丛} $\mathbf V(N)$ 的\Term{projective-completion}{射影完备化}。与 9.1 的陈述相同，以 $s:Y\hookrightarrow\widehat N$ 表示 $\mathbf V(N)$ 的\Term{zero-section}{零截面}与\Term{canonical}{典范}\Term{inclusion}{包含映射} $\mathbf V(N)\hookrightarrow\widehat N$ 的复合。

\Term{inclusion}{包含映射} $Y\subset X$ 使我们可将 $Y\times\mathbf P^1=W$ 等同于 $X\times\mathbf P^1$ 的一个\Term{closed-subscheme}{闭子概形}。以 $W'$ 表示“$Y\times\mathbf P^1$ 在 $f_1$ 下的\Term{strict-transform}{严格变换}（transformé pur）”，即 $Z'$ 的\Term{closed-subscheme}{闭子概形}，它是 $Y\times\mathbf P^1$ 沿 $Y=Y\times0$ 的\Term{blow-up}{爆破}。显然 $W'\cap\widehat N$ 可等同于 $\widehat N$ 的\Term{zero-section}{零截面}；且由于 $Y$ 在 $Y\times\mathbf P^1$ 中的\Term{codimension}{余维数}为 $1$，\Term{canonical}{典范}\Term{morphism}{态射} $f_2:W'\to Y\times\mathbf P^1$ 为\Term{isomorphism}{同构}。关于 $W'$ 的其他记号由下列\Term{canonical}{典范}图确定：
\SourcePage{350}
\[
\begin{tikzcd}[column sep=large,row sep=large]
 &\widehat N\arrow[r,hook,"j_1"]\arrow[ddl,"g_1",bend right=12]&Z'\arrow[dd,"f_1"]\\
Y\arrow[ur,hook,"s"]\arrow[r,hook,"\delta"]\arrow[d,"\mathrm{id}"']&W'\arrow[ur,hook,"\beta"]\arrow[d,"f_2","\wr"']&\\
Y\arrow[r,hook,"\gamma"]&Y\times\mathbf P^1\arrow[r,hook,"\alpha"]&X\times\mathbf P^1
\end{tikzcd}\tag{$D_1$}
\]
同样，以 $X'$ 表示 $Z'$ 的\Term{closed-subscheme}{闭子概形}，它是 $X=X\times0$ 沿 $Y$ 的\Term{blow-up}{爆破}。显然 $\mathbf P(N)=X'\cap\widehat N$ 可等同于 $\widehat N$ 的\Term{hyperplane-at-infinity}{无穷远超平面} $H$。关于 $X'$ 的其他记号由下列\Term{canonical}{典范}图确定：
\[
\begin{tikzcd}[column sep=large,row sep=large]
 &\widehat N\arrow[r,hook,"j_1"]\arrow[ddl,"g_1",bend right=12]&Z'\arrow[dd,"f_1"]\\
H\arrow[ur,hook,"k"]\arrow[r,hook,"j"]\arrow[d,"g"']&X'\arrow[ur,hook,"v"]\arrow[d,"f"]&\\
Y\arrow[r,hook,"i"]&X\arrow[r,hook,"u"]&X\times\mathbf P^1=Z
\end{tikzcd}\tag{$D_2$}
\]
下面以\Term{lemma}{引理}的\Term{form}{形式}给出若干后文所需的关系。以 $E$ 表示由\Term{canonical}{典范}\Term{exact-sequence}{正合列}
\[
 0\longrightarrow E\longrightarrow g_1^*(N\oplus\mathscr O_Y)\longrightarrow\mathscr O_{\widehat N}(1)\longrightarrow0
 \tag{9.3.0}
\]
定义的\Term{locally-free}{局部自由} $\mathscr O_{\widehat N}$-\Term{module}{模}，
\SourcePage{351}
并置
\[
 \bar\xi=c_1(\mathscr O_{\widehat N}(1)).\tag{9.3.1}
\]
\begin{SGAresult}{\Term{lemma}{引理}}{9.3}
沿用上述记号，有关系
\[
 j_1^*(\beta_*(1))=c_d(\check E),\quad\text{在 }C^\bullet(\widehat N)\text{ 中},\tag{9.3.2}
\]
\[
 j_1^*(v_*(1))=\bar\xi,\quad\text{在 }C^\bullet(\widehat N)\text{ 中},\tag{9.3.3}
\]
\[
 f_1^*u_*(1)=(j_1)_*(1)+v_*(1),\quad\text{在 }C^\bullet(Z')\text{ 中}.\tag{9.3.4}
\]
\end{SGAresult}
\begin{proof}
所考虑的 $k$-\Term{scheme}{概形}均\Term{smooth}{光滑}，故由\Term{dimension}{维数}立即可知\Term{morphism}{态射} $j_1$ 和 $\beta$（分别地，$j_1$ 和 $v$）\Term{transverse}{横截}。由 [CH] 4 命题 2 的推论 1 得等式
\[
 j_1^*\beta_*(1)=s_*(1),\qquad j_1^*v_*(1)=k_*(1).
\]
等式 (9.3.2) 和 (9.3.3) 分别由 (9.1.2) 及\Term{chern-class}{陈类}的\Term{normalization}{正规化}（[CH] 4--18）得到。再证 (9.3.4)。$Z'$ 中的\Term{cycle}{圈} $X'$ 是 $X\times0$ 在 $f_1$ 下的\Term{strict-transform}{严格变换}，故先验地存在形如
\[
 f_1^*u_*(1)=(j_1)_*g_1^*(\mu)+v_*(1),\tag{9.3.5}
\]
的等式，其中 $\mu\in\Gamma(Y,\mathbf N)$；须证 $\mu=1$。这可通过改写 SGA 7 第 X 讲 2.6.3 的证明做到。这里给出直接证明。对 (9.3.5) 两边施以 $j_1^*$，得（参见 $D_2$）
\[
 g_1^*i^*u^*u_*(1)=j_1^*(j_1)_*g_1^*(\mu)+j_1^*v_*(1).\tag{9.3.6}
\]
(9.1.3) 的一个极特殊情形，或直接验证，表明 $u^*u_*(1)=0$。再用 (9.3.3)，由 (9.3.6) 得
\[
 0=j_1^*(j_1)_*g_1^*(\mu)+\bar\xi.\tag{9.3.7}
\]
\SourcePage{352}
取 $Y$ 的一个\Term{connected-component}{连通分支}的\Term{neighborhood}{邻域}后，可设 $Y$ \Term{connected}{连通}，于是 (9.3.7) 化为
\[
 0=\mu\,j_1^*(j_1)_*(1)+\bar\xi,\qquad\mu\in\mathbf N.\tag{9.3.7 bis}
\]
因为 $j_1^*(j_1)_*(1)$ 是 $1$ 次\Term{homogeneous}{齐次}元素，它具有\Term{form}{形式}
\[
 j_1^*(j_1)_*(1)=g_1^*(\rho)+\nu\bar\xi,
 \qquad\rho\in C^1(Y),\quad\nu\in\mathbf Z.
\]
代入 (9.3.7 bis)，立即得 $\rho=0$ 及 $\mu\nu=-1$，而 $\mu\in\mathbf N$，故 $\mu=1$、$\nu=-1$；这就完成了 (9.3.4) 的证明。
\end{proof}
9.2 的证明关键在于证明下一个\Term{lemma}{引理}（若像\Term{cohomology}{上同调}中那样，已知对每个 $y\in C^\bullet(\widehat N)$ 有 $j_1^*(j_1)_*(y)=-y\bar\xi$，则它显然成立）。
\begin{SGAresult}{\Term{lemma}{引理}}{9.4}
对每个 $y\in C^\bullet(\widehat N)$，有
\[
 c_d(\check E)j_1^*(j_1)_*(y)=0.
\]
\end{SGAresult}
\begin{proof}
元素 $y$ 在 $C^\bullet(\widehat N)$ 中有唯一\Term{resolution}{分解}
\[
 y=\sum_{0\le p\le d}g_1^*(a_p)\bar\xi^p,
\]
即由 (9.3.3)，
\[
 y=\sum_{0\le p\le d}g_1^*(a_p)[j_1^*v_*(1)]^p.\tag{9.4.1}
\]
由\Term{morphism}{态射} $j_1$ 的\Term{projection-formula}{投影公式}（[CH] 4 I 9）及 (9.1.6)，得到
\[
 c_d(\check E)j_1^*(j_1)_*(y)=c_d(\check E)j_1^*(j_1)_*g_1^*(a_0).
\]
换言之，只需在 $y$ 形如 $g_1^*(a)$、$a\in C^\bullet(Y)$ 时证明 9.4。此时元素 $j_1^*(j_1)_*(g_1)^*(a)$ 在 $C^\bullet(\widehat N)$ 中又有唯一\Term{resolution}{分解}
\SourcePage{353}
\[
 j_1^*(j_1)_*g_1^*(a)=\sum_{0\le p\le d}g_1^*(b_p)\bar\xi^p.\tag{9.4.2}
\]
再次利用 (9.1.6)，须证 $g_1^*(b_0)c_d(\check E)=0$。事实上将证明更强的 $b_0=0$（依 9.1.1，这两者\Term{equivalence}{等价}）。由 (9.3.3)，可将 (9.4.2) 改写为
\[
 j_1^*(j_1)_*g_1^*(a)=g_1^*(b_0)+z\,j_1^*v_*(1),\quad z\in C^\bullet(\widehat N).\tag{9.4.3}
\]
对 (9.4.3) 两边施以 $(j_1)_*$，由 $j_1$ 的\Term{projection-formula}{投影公式}得
\[
 (j_1)_*g_1^*(b_0)=[(j_1)_*g_1^*(a)]\cdot(j_1)_*(1)-v_*(1)\cdot(j_1)_*(z).
 \tag{9.4.4}
\]
将由 (9.3.4) 得到的 $(j_1)_*(1)$ 的表达式代入 (9.4.4)，得
\[
 (j_1)_*g_1^*(b_0)=(j_1)_*g_1^*(a)\cdot f_1^*u_*(1)
 -v_*(1)[(j_1)_*g_1^*(a)+(j_1)_*(z)].\tag{9.4.5}
\]
由于 $X\times0$ 与 $X\times1$ 在 $X\times\mathbf P^1$ 中\Term{linear-equivalence}{线性等价}，有等式 $f_1^*u_*(1)=f_1^*(X\times1)$；而 $\widehat N\cap f_1^{-1}(X\times1)=\varnothing$，故
\[
 (j_1)_*g_1^*(a)\cdot f_1^*u_*(1)=0.\tag{9.4.6}
\]
与 (9.4.5) 比较，得到
\[
 (j_1)_*g_1^*(b_0)=-v_*(1)[(j_1)_*g_1^*(a)+(j_1)_*(z)].\tag{9.4.7}
\]
另一方面，显然 $W'\cap X'=\varnothing$：它们至多在 $Y$ 上方相交，而在那里一个给出 $\widehat N$ 的\Term{zero-section}{零截面}，另一个给出\Term{hyperplane-at-infinity}{无穷远超平面}。因而 $\beta^*v_*(1)=0$，由 (9.4.7) 得
\[
 \beta^*(j_1)_*g_1^*(b_0)=0.\tag{9.4.8}
\]
但如前所述（参见 9.3.2 的证明），$\beta$ 与 $j_1$ \Term{transverse}{横截}，所以 (9.4.8) 给出
\[
 0=\delta_*s^*g_1^*(b_0)=\delta_*(b_0).\tag{9.4.9}
\]
\SourcePage{354}
最后，对 (9.4.9) 施以 $(f_2)_*$，得到 $\gamma_*(b_0)=0$，故 $b_0=0$，因为 $\gamma_*$ 为\Term{injective}{单射}（$\gamma$ 显然有一个\Term{retraction}{收缩映射}）。
\end{proof}
\textbf{9.5。}最后说明上述结果如何\Term{pushout}{推出}\Term{theorem}{定理}。以 $p:X\times\mathbf P_k^1\to X$ 表示第一\Term{projection}{投影}，于是 $i=p\circ i_1$。若 $y\in C^\bullet(Y)$，显然有
\[
 i^*i_*(y)=(i_1)^*(\bar x),\qquad\bar x=p^*i_*(y).\tag{9.5.1}
\]
由 (9.1.1) 得 $i^*i_*(y)=(g_1)_*[g_1^*i_1^*(\bar x)c_d(\check E)]$，即
\[
 i^*i_*(y)=(g_1)_*[j_1^*f_1^*(\bar x)\cdot c_d(\check E)].\tag{9.5.2}
\]
为证明\Term{theorem}{定理}，由\Term{linear}{线性}可设 $y$ 为 $Y$ 的\Term{subvariety}{子簇} $T$ 对应的\Term{irreducible-cycle}{不可约圈}，于是 $\bar x$ 对应于 $Z=X\times\mathbf P^1$ 中的\Term{cycle}{圈} $T\times\mathbf P^1$。以 $\theta$ 表示 $T\times\mathbf P^1$ 在 $f_1$ 下的\Term{strict-transform}{严格变换}，即 $Z\setminus Y$ 上方 $T\times\mathbf P^1$ 的\Term{inverse-image}{逆像}的\Term{closure}{闭包}，则存在形如
\[
 f_1^*(\bar x)=\theta+(j_1)_*(w),\qquad w\in C^\bullet(\widehat N)\tag{9.5.3}
\]
的关系。将此式代入 (9.5.2)，由 9.4 得
\[
 i^*i_*(y)=(g_1)_*[j_1^*(\theta)c_d(\check E)],\tag{9.5.4}
\]
再由 (9.1.4)，
\[
 i^*i_*(y)=s^*j_1^*(\theta).\tag{9.5.5}
\]
而 $\theta$ 的底层\Term{reduced-scheme}{约化概形}是 $T\times\mathbf P^1$ 沿 $T=(T\times\mathbf P^1)\cap Y$ 的\Term{blow-up}{爆破}，故它与 $\widehat N$ 的\Term{zero-section}{零截面}\Term{transverse}{横截}相交，交集正是 $T$ 在该\Term{zero-section}{零截面}下的像。因此 $j_1^*(\theta)=s_*(y)$，结合 (9.5.5) 和 (9.1.3)，即完成\Term{theorem}{定理}的证明。
\SourcePage{355}
现在转向导言所说的“关键公式”。所用记号主要沿用第 4 节，并在此大多重述。注意这些记号与本节引入的记号相容，唯一的区别是第 4 节中记为 $Y'$ 的\Term{scheme}{概形}在这里记作 $H$（取自“hyperplan”，即\Term{hyperplane}{超平面}）。下面也将交替使用这两个记号。
\begin{SGAresult}{命题}{9.6}
设 $k$ 为\Term{separably-closed}{可分闭}\Term{field}{域}，$i:Y\hookrightarrow X$ 为\Term{quasi-projective}{拟射影}\Term{smooth}{光滑} $k$-\Term{scheme}{概形}间的\Term{closed-immersion}{闭浸入}，纯\Term{codimension}{余维数}为 $d$。以 $f:X'\to X$ 表示将 $X$ 沿 $Y$ \Term{blow-up}{爆破}所得的 $X$-\Term{scheme}{概形}，设
\[
\begin{tikzcd}
Y'\arrow[r,hook,"j"]\arrow[d,"g"']&X'\arrow[d,"f"]\\
Y\arrow[r,hook,"i"]&X
\end{tikzcd}
\]
是由 $f$ 与 $i$ 构造的\Term{cartesian-diagram}{笛卡尔图}。以 $F$ 表示由\Term{canonical}{典范}\Term{exact-sequence}{正合列}
\[
 0\longrightarrow F\longrightarrow g^*(N)\longrightarrow\mathscr O_{Y'}(1)\longrightarrow0
 \tag{9.6.1}
\]
定义的\Term{locally-free}{局部自由} $\mathscr O_{Y'}$-\Term{module}{模}，则对每个 $y\in C^\bullet(Y)$，在 $C^\bullet(X')$ 中有等式
\[
 f^*i_*(y)=j_*(g^*(y)c_{d-1}(\check F)).\tag{9.6.2}
\]
\end{SGAresult}
在证明 9.6 之前，先将图 $D_1$、$D_2$ 的一些尚未证明的性质整理成\Term{lemma}{引理}。
\begin{SGAresult}{\Term{lemma}{引理}}{9.7}
有等式：
\[
\begin{aligned}
\text{(i)}\quad c_d(\check E)&=k_*(c_{d-1}(\check F))+(g_1)^*(c_d(\check N));\\
\text{(ii)}\quad\beta_*f_2^*\gamma_*(y)&=f_1^*\alpha_*\gamma_*(y),\quad y\in C^\bullet(Y);\\
\text{(iii)}\quad f_1^*u_*(t)&=(j_1)_*[(i\circ g_1)^*(t)]+v_*f^*(t),\quad t\in C^\bullet(X).
\end{aligned}
\]
\end{SGAresult}
\SourcePage{356}
\begin{proof}
先证 (i)。置 $\xi=c_1(\mathscr O_{Y'}(1))=k^*(\bar\xi)$。仿照 (8.4.1) 的证明得到等式
\[
 c_d(\check E)=\sum_{p=0}^d g_1^*(c_p(\check N))\bar\xi^{d-p},\tag{9.7.1}
\]
\[
 c_{d-1}(\check F)=\sum_{q=0}^{d-1}g^*(c_q(\check N))\xi^{d-1-q}.\tag{9.7.2}
\]
由 (9.7.2) 得
\[
 k_*[c_{d-1}(\check F)]=\sum_{q=0}^{d-1}k_*k^*[g_1^*(c_q(\check N))\bar\xi^{d-1-q}],
\]
再用 $k$ 的\Term{projection-formula}{投影公式}及 $k_*(1)=\bar\xi$，得
\[
 k_*[c_{d-1}(\check F)]=\sum_{q=0}^{d-1}g_1^*(c_q(\check N))\bar\xi^{d-q}.\tag{9.7.3}
\]
比较 (9.7.1) 与 (9.7.3)，立即得到 (i)。再证 (ii)。可设 $y$ 为 $Y$ 的\Term{irreducible-cycle}{不可约圈}，对应的\Term{subvariety}{子簇}也记为 $y$。于是\Term{cycle}{圈} $\gamma_*(y)=y\times0$ 与 $y\times1$ \Term{rational-equivalence}{有理等价}，故只须证明\Term{cycle}{圈}的等式
\[
 \beta_*f_2^*(y\times1)=f_1^*\alpha_*(y\times1).\tag{9.7.4}
\]
但 $(y\times1)\cap(Y\times0)=\varnothing$，故 $f_1$ 在 $y\times1$ 的一个\Term{open-neighborhood}{开邻域}上方为\Term{isomorphism}{同构}，立即得到 (9.7.4)。再证 (iii)。以 $p:X\times\mathbf P^1\to X$ 表示第一\Term{projection}{投影}，由 (9.1.2) 得
\[
 u_*(t)=p^*(t)\cdot u_*(1),\tag{9.7.5}
\]
与 (9.3.4) 比较，得
\[
 f_1^*u_*(t)=(p\circ f_1)^*(t)[(j_1)_*(1)+v_*(1)].\tag{9.7.6}
\]
分别利用 $j_1$ 与 $v$ 的\Term{projection-formula}{投影公式}，得
\[
 f_1^*u_*(t)=(j_1)_*[(p\circ f_1\circ j_1)^*(t)]+v_*[(p\circ f_1\circ v)^*(t)].
 \tag{9.7.7}
\]
因 $p\circ f_1\circ j_1=(p\circ u)\circ i\circ g_1=i\circ g_1$，且 $p\circ f_1\circ v=(p\circ u)\circ f=f$，便得 (iii)。
\end{proof}
\SourcePage{357}
\textbf{9.8。}现在证明 9.6。为此先证等式
\[
 v_*f^*i_*(y)=v_*j_*(g^*(y)c_{d-1}(\check F)),\tag{9.8.1}
\]
然后证明 $v_*$ 为\Term{injective}{单射}。由 (9.7) (iii) 得
\[
 v_*f^*i_*(y)=f_1^*u_*i_*(y)-(j_1)_*[g_1^*i^*i_*(y)],
\]
由 (9.2) 即
\[
 v_*f^*i_*(y)=f_1^*(i_1)_*(y)-(j_1)_*g_1^*(y\,c_d(\check N)).\tag{9.8.2}
\]
另一方面，立即有
\[
 v_*j_*(g^*(y)c_{d-1}(\check F))=(j_1)_*[(k_*k^*g_1^*(y))c_{d-1}(\check F)],
\]
即由 $k$ 的\Term{projection-formula}{投影公式}，
\[
 v_*j_*(g^*(y)c_{d-1}(\check F))=(j_1)_*[g_1^*(y)k_*(c_{d-1}(\check F))].
\]
利用 (9.7) (i) \Term{pushout}{推出}
\[
\begin{aligned}
 v_*j_*(g^*(y)c_{d-1}(\check F))
 &=(j_1)_*(g_1^*(y)c_d(\check E))\\
 &\quad-(j_1)_*g_1^*(y\,c_d(\check N)).
\end{aligned}\tag{9.8.3}
\]
比较 (9.8.2) 与 (9.8.3)，可见 (9.8.1) \Term{equivalence}{等价}于
\[
 f_1^*(i_1)_*(y)=(j_1)_*(g_1^*(y)c_d(\check E)),\tag{9.8.4}
\]
即将\SGAsourcegap{原文在“lorsqu'on remplace par”中未写出被替换的\Term{object}{对象}}换为 $X\times\mathbf P^1$ 时，与 (9.6.2) 对应的等式。由 (9.1.2)，有
\[
 (j_1)_*(g_1^*(y)c_d(\check E))=(j_1)_*s_*(y)=\beta_*\delta_*(y).
\]
又显然（参见 $D_1$）$\delta_*=f_2^*\gamma_*$，故
\[
 (j_1)_*(g_1^*(y)c_d(\check E))=\beta_*f_2^*\gamma_*(y).\tag{9.8.5}
\]
另一方面，显然
\[
 f_1^*(i_1)_*(y)=f_1^*\alpha_*\gamma_*(y),\tag{9.8.6}
\]
所以比较 (9.8.5)、(9.8.6) 及 (9.7) (ii)，即得 (9.8.4)。
\SourcePage{358}
\begin{SGAresult}{\Term{lemma}{引理}}{9.8.7}
映射 $v_*$ 为\Term{injective}{单射}。
\end{SGAresult}
\begin{proof}
设 $Z''$ 为 $Z'$ 沿 $W'$ 的\Term{blow-up}{爆破}，\Term{equivalence}{等价}地，是沿 $\widehat N\cup W'=f_1^{-1}(Y\times\mathbf P^1)$ 的\Term{blow-up}{爆破}。$X$ 在 $Z''$ 中的\Term{strict-transform}{严格变换} $X''$ 也是 $X'$ 在 $Z''$ 中的\Term{strict-transform}{严格变换}。记号由以下\Term{canonical}{典范}图确定：
\[
\begin{tikzcd}
X''\arrow[r,hook,"v''"]\arrow[d,"f'","\wr"']&Z''\arrow[d,"f_1'"]\\
X'\arrow[r,hook,"v"]&Z'
\end{tikzcd}
\]
由于 $X'\cap W'=\varnothing$，\Term{morphism}{态射} $f_1'$ 在 $X'$ 的一个\Term{open-neighborhood}{开邻域}上方为\Term{isomorphism}{同构}，特别地，$f'$ 为\Term{isomorphism}{同构}。又因为 $f_1'$ 与 $v$ \Term{transverse}{横截}，有等式
\[
 (f_1')^*\circ v_*=v_*''\circ(f')^*,\tag{9.8.8}
\]
故只须证明 $v_*''$ 为\Term{injective}{单射}。但由\Term{blow-up}{爆破}的\Term{commutativity}{交换性}（参见例如 Hironaka，《\Term{algebraic-cycle}{代数圈}的\Term{smooth}{光滑}化》（Smoothing of algebraic cycles），\emph{Amer. Journal of Math.}（1968），\Term{lemma}{引理} 4.1），$Z''$ 也可先将 $X\times\mathbf P^1$ 沿 $Y\times\mathbf P^1$ \Term{blow-up}{爆破}得到\Term{scheme}{概形} $Z_1'$，再将 $Z_1'$ 沿 $Y$ 的\Term{inverse-image}{逆像}\Term{blow-up}{爆破}而得。显然 $Z_1'=X'\times\mathbf P^1$，其\Term{projection}{投影} $Z_1'\to X\times\mathbf P^1$ 为乘积 $f\times\mathrm{id}_{\mathbf P^1}$，由此有显然的\Term{cartesian-diagram}{笛卡尔图}
\[
\begin{tikzcd}[column sep=large]
H\arrow[r,hook,"{t\mapsto(t,0)}"]\arrow[d,"g"']&H\times\mathbf P^1\arrow[r,hook,"j\times\mathrm{id}"]\arrow[d,"g\times\mathrm{id}"]&X'\times\mathbf P^1=Z_1'\arrow[d,"f\times\mathrm{id}"]\\
Y\arrow[r,hook,"\gamma"]&Y\times\mathbf P^1\arrow[r,hook,"i\times\mathrm{id}"]&X\times\mathbf P^1
\end{tikzcd}
\]
$X\times0$ 在 $f\times\mathrm{id}$ 下的\Term{strict-transform}{严格变换}可等同于 $X'\times0$，它包含 $H=H\times0=(f\times\mathrm{id})^{-1}(Y)$。现在将 $Z_1'$ 沿 $H$ \Term{blow-up}{爆破}，由于 $H$ 在 $X'$ 中纯\Term{codimension}{余维数}为 $1$，$X'$ 的\Term{strict-transform}{严格变换}与 $X'$ \Term{canonical}{典范}\Term{isomorphic}{同构}（这一点也已知），故有\Term{canonical}{典范}\Term{commutative-diagram}{交换图}
\SourcePage{359}
\[
\begin{tikzcd}[column sep=large]
 &X'\arrow[r,hook,"v''"]\arrow[d,"f''"]&Z''\arrow[d,"f_1''"]\\
 H\arrow[r,hook]&X'\arrow[r,hook,"v'"']&X'\times\mathbf P^1\\
 &t\arrow[r,mapsto]&(t,0)
\end{tikzcd}
\]
回忆要证明 $v_*''$ 为\Term{injective}{单射}。与 $D_2$ 比较可见，这\Term{equivalence}{等价}于证明在额外假定 $Y$ 在 $X$ 中的\Term{codimension}{余维数}为 $1$ 时，$v_*$ 为\Term{injective}{单射}；下面即证明这一点。此时等式 $(f_1)_*\circ v_*=u_*\circ f_*$ 即给出结论，因为 $f_*$ 是\Term{isomorphism}{同构}，$u_*$ 是\Term{split-monomorphism}{分裂单态射}（$u$ 有\Term{retraction}{收缩映射}）。
\end{proof}
现在陈述与 (8.5) b) 对应的\Term{theorem}{定理}。
\begin{SGAresult}{\Term{theorem}{定理}}{9.9}
对每个整数 $p\ge0$，序列
\[
 0\longrightarrow C^{p-d}(Y)\xrightarrow{\lambda}C^{p-1}(Y')\oplus C^p(X)
 \xrightarrow{\mu}C^p(X')\longrightarrow0,
\]
其中 $\lambda$、$\mu$ 表示\Term{arrow}{箭头}
\[
 \begin{aligned}
 \lambda:x&\longmapsto[g^*(x)c_{d-1}(\check F),-i_*(x)],\\
 \mu:(x,y)&\longmapsto j_*(x)+f^*(y),
 \end{aligned}
\]
是\Term{exact-sequence}{正合列}。此外，映射 $\lambda$ 有\Term{left-inverse}{左逆}
\[
 \lambda':(x,y)\longmapsto g_*(x).
\]
\end{SGAresult}
\begin{proof}
最后一条断言正是 (9.1.1)。同样，(9.6.2) 表示 $\mu\circ\lambda=0$。现证 $\mu$ 为\Term{surjective}{满射}。为此需用到
\[
 f_*f^*(z)=z\qquad(z\in C^\bullet(X)),\tag{9.9.1}
\]
它直接由 $f$ 的\Term{projection-formula}{投影公式}及显然的事实 $f_*(1)=1$ 得到。设 $t\in C^\bullet(X')$。由 (9.9.1)，有 $f_*(t-f^*f_*(t))=0$。由于 $f$ 诱导 $U'=X'\setminus Y'$ 到 $U=X\setminus Y$ 的\Term{isomorphism}{同构}，可知
\SourcePage{360}
$t-f^*f_*(t)$ 在 $U'$ 上的\Term{restriction}{限制}为零，故由 [CH] 4 \Term{theorem}{定理} 2，有等式
\[
 t-f^*f_*(t)=j_*(a),\qquad a\in C^\bullet(Y').
\]
现证中间的\Term{exactness}{正合性}。设 $a\in C^{p-1}(Y')$、$b\in C^p(X)$ 满足
\[
 j_*(a)+f^*(b)=0.\tag{9.9.2}
\]
对 (9.9.2) 两边施以 $f_*$，得到 $f_*j_*(a)+f_*f^*(b)=0$，再由 (9.9.1)，
\[
 b=-i_*g_*(a).\tag{9.9.3}
\]
置 $x=g_*(a)$，则只须证明
\[
 a=g^*(x)c_{d-1}(\check F).\tag{9.9.4}
\]
为此置 $z=a-g^*(x)c_{d-1}(\check F)$（$x=g_*(a)$），并证明 $z=0$。先证明有关系 $j_*(z)=0$、$g_*(z)=0$，再证明满足这两个关系的元素 $z$ 必为零。依定义，
\[
 j_*(z)=j_*(a)-j_*(g^*(x)\cdot c_{d-1}(\check F)),
\]
由 (9.6.2) 即 $j_*(z)=j_*(a)-f^*i_*(x)$。在此式中将 $i_*(x)$ 换为 $-b$（9.9.3），便由 (9.9.2) 得 $j_*(z)=0$。再由 $x$ 的定义及\Term{theorem}{定理}最后一条断言，得
\[
 g_*(z)=x-g_*(g^*(x)c_{d-1}(\check F))=0.
\]
下面说明如何由 $j_*(z)=0$ 与 $g_*(z)=0$ 两个关系\Term{pushout}{推出}
\SourcePage{361}
$z=0$。由于 $Y'\simeq\mathbf P_Y(N)$，元素 $z$ 有唯一表示
\[
 z=g^*(a_0)+g^*(a_1)\xi+\cdots+g^*(a_{d-1})\xi^{d-1},
 \qquad a_i\in C^{p-1-i}(Y).
\]
由于 $g_*(z)=0$，由 $g$ 的\Term{projection-formula}{投影公式}以及 $g_*$ 将\Term{degree}{次数}降低 $d-1$ 这一事实，得 $a_{d-1}=0$。另一方面，由 (9.2)，有 $z\xi=j^*j_*(z)=0$，故
\[
 g^*(a_0)\xi+\cdots+g^*(a_{d-1})\xi^{d-1}=0;
\]
所以所有 $a_i$ 都为零，$z=0$。
\end{proof}
\textbf{9.10。}沿用 (9.9) 的记号，映射 $\omega=\lambda\circ\lambda'$ 是一个\Term{projector}{投影算子}，满足 $\operatorname{Im}(\omega)=\operatorname{Im}(\lambda)$。因而 $\mu$ 诱导\Term{isomorphism}{同构}
\[
 \bar\mu:\operatorname{Ker}(\omega)\xrightarrow{\sim}C^p(X').
\]
而
\[
 \operatorname{Ker}(\omega)=\operatorname{Ker}(g_*)\oplus C^p(X)
 \hookrightarrow C^{p-1}(Y')\oplus C^p(X).\tag{9.10.1}
\]
事实上，$C^{p-1}(Y')\oplus C^p(X)$ 中的元素 $(x,y)$ 属于 $\operatorname{Ker}(\omega)$，当且仅当
\[
 g^*g_*(x)\cdot c_{d-1}(\check F)=-i_*g_*(x)=0,
\]
亦即由 (9.1.1)，$g_*(x)=0$。

另一方面，\Term{projective-bundle}{射影丛}的\Term{chow-ring}{Chow 环}结构，尤其是关系
\[
 g_*(\xi^i)=0\quad(i<d-1),\qquad g_*(\xi^{d-1})=1,
\]
表明
\[
 \operatorname{Ker}(g_*)=g^*C^{p-1}(Y)\oplus g^*C^{p-2}(Y)\xi\oplus\cdots\oplus g^*C^{p-d+1}(Y)\xi^{d-2}.
\]
换言之，$\bar\mu$ \Term{solvable}{可解}释为\Term{isomorphism}{同构}
\SourcePage{362}
\[
 C^{p-1}(Y)\oplus C^{p-2}(Y)\oplus\cdots\oplus C^{p-d+1}(Y)\oplus C^p(X)
 \xrightarrow{\sim}C^p(X'),\tag{9.10.2}
\]
其逆\Term{isomorphism}{同构}也容易写出。

正如 Murre 告诉我的，Samuel 已在 $d=2$ 时证明了命题 (9.10.2)（《\Term{equivalence-relation}{等价关系}》（Relations d'équivalence），\emph{Proceedings of the International Congress}，Edinburgh 1958）。

当然，在\Term{etale-cohomology}{平展上同调}或\Term{l-adic-cohomology}{$\ell$-进上同调}的框架内，也有与 (9.10.2) 对应的命题，这时可由 (8.5) \Term{pushout}{推出}。

\textbf{参考文献}
\begin{description}
\item[{[1]}] Grothendieck，\emph{\Term{chern-class}{陈类}理论}（La théorie des classes de Chern）。\emph{Bull. Soc. Math. de France} \textbf{86}（1958），第 137--154 页。
\item[{[2]}] Hodge，\emph{调和积分的理论及应用}（The theory and applications of harmonic integrals）。Cambridge University Press，1952。
\item[{[3]}] Hirzebruch，\emph{\Term{algebraic-geometry}{代数几何}中的\Term{topology}{拓扑}方法}（Topological Methods in Algebraic Geometry）。Springer Verlag。
\item[{[4]}] Lefschetz，\emph{\Term{algebraic-geometry}{代数几何}}（Algebraic Geometry）。Princeton University Press，1953。
\item[{[5]}] Raynaud 女士，\emph{泛\Term{projective}{射影}模}（Modules projectifs universels）。\emph{Inventiones Mathematicae}，第 6 卷，第 1 分册（1968），第 1--26 页。
\item[{[TF]}] Godement，\emph{层论}（Théorie des faisceaux）。
\item[{[CH]}] \emph{Chow 环及其应用}（Anneaux de Chow et applications）。Séminaire Chevalley，1958。
\item[{[6]}] \emph{Lefschetz \Term{theorem}{定理}及\Term{degeneration}{退化}判据}（Théorème de Lefschetz et critères de dégénérescence）。\emph{Pub. I.H.E.S.}，第 35 号。
\end{description}
